\newpage
\section{Przegląd literatury i analiza rozwiązań}

\subsection{Architektury hybrydowe Edge--Cloud}


Wdrożenie asystenta głosowego działającego lokalnie wymaga rozstrzygnięcia między dwiema skrajnymi architekturami: przetwarzaniem w chmurze (ang. \emph{cloud computing}) oraz przetwarzaniem na urządzeniu brzegowym (ang. \emph{edge computing}). Każda z nich niesie odmienne kompromisy, które przekładają się bezpośrednio na jakość i użyteczność systemu konwersacyjnego.

Architektura chmurowa zapewnia dostęp do dużych modeli oraz praktycznie nieograniczonej, skalowalnej mocy obliczeniowej, lecz obarczona jest opóźnieniem wynikającym z transmisji sieciowej, które ma znaczenie krytyczne dla naturalnej konwersacji, a ponadto uzależnia działanie systemu od stałego połączenia z internetem. W środowiskach szczególnie wrażliwych na kwestie prywatności, takich jak obiekty rekreacyjne czy inteligentne domy, przesyłanie surowych danych pochodzących z czujników do zewnętrznych usług budzi uzasadnione obawy \cite{Renney2026}.

Architektura brzegowa umożliwia z kolei przetwarzanie lokalne, pozbawione opóźnień sieciowych i niewymagające udostępniania danych podmiotom trzecim, jednak istotnie ogranicza dopuszczalny rozmiar wykorzystywanych modeli. Badania empiryczne wskazują, że procesory ARM CPU pozbawione akceleratora GPU osiągają wydajność rzędu 5--15 tokenów na sekundę dla modeli o wielkości 1,5~mld parametrów oraz 2--5 tokenów na sekundę dla modeli 3~mld parametrów \cite{Nguyen2025SBC}, przy czym dominującym ograniczeniem pozostaje przepustowość pamięci (ang. \emph{memory bandwidth}) oraz konieczność zarządzania temperaturą procesora \cite{Tummalapalli2026}.

Rozwiązaniem kompromisowym jest architektura hybrydowa, w której ciężkie rozumowanie (ang. \emph{reasoning}) realizowane jest na lokalnym serwerze, natomiast inferencja bezpośrednio obsługująca użytkownika odbywa się na urządzeniu brzegowym. Podejście to zachowuje prywatność danych, jednocześnie zapewniając dostęp do modeli wyższej jakości na potrzeby złożonej analizy \cite{Renney2026}. Architektura Wilga realizuje ten wzorzec poprzez podział ról Nauczyciel--Uczeń, opisany szczegółowo w sekcji~\ref{subsec:teacher_student}.


Fundamentem współczesnego przetwarzania języka naturalnego pozostaje architektura Transformer \cite{Vaswani2017}, na której opierają się niemal wszystkie stosowane obecnie modele językowe. W ostatnich latach wyraźnie zaznacza się trend zwrotu ku małym modelom językowym (ang. \emph{Small Language Models}, SLM), liczącym poniżej 8~mld parametrów, które dzięki treningowi prowadzonemu na wysokiej jakości, selektywnie dobranych danych dorównują jednostkom wielokrotnie od nich większym. Potwierdzają to raporty techniczne opublikowane przez firmy Meta \cite{Llama3}, Microsoft \cite{Phi3} oraz Google \cite{Gemma2}: modele Llama~3, Phi-3 i Gemma~2 osiągają wyniki konkurencyjne wobec modeli rzędu dziesiątek miliardów parametrów, przy ułamku ich wymagań pamięciowych. Uzupełnieniem tego nurtu jest Bielik \cite{Bielik}, polski model oparty na architekturze Mistral, rozwijany w ramach konsorcjum SpeakLeash.

Pole małych modeli językowych rozwija się w sposób wyjątkowo dynamiczny. Już w trakcie realizacji niniejszej pracy pojawiły się kolejne generacje modeli, w tym Gemma~3 oraz Gemma~4, co wymuszało bieżące korygowanie decyzji dotyczących zakresu ewaluacji. Renney i in. \cite{Renney2026} wskazują na wysoką wymiarowość przestrzeni konfiguracji jako istotne utrudnienie w badaniach nad inferencją brzegową: konfiguracja optymalna dla jednej generacji modeli może nie być transferowalna na generację następną. Z tego względu zastosowana w pracy metodologia benchmarku, oparta na automatycznej ewaluacji typu LLM-as-judge oraz kompozytowym wskaźniku jakości, została zaprojektowana jako niezależna od konkretnego modelu (ang. \emph{model-agnostic}) i powtarzalna bez konieczności modyfikacji infrastruktury.

Dla platformy Raspberry Pi~5, wyposażonej w 8~GB pamięci RAM oraz procesor ARM Cortex-A76, klasa modeli o wielkości od 2 do 4~mld parametrów stanowi praktyczne optimum. Modele mniejsze niż 1,5~mld parametrów często nie osiągają zadowalającej jakości rozumowania w języku polskim, natomiast modele przekraczające 4~mld parametrów wyczerpują dostępną pamięć operacyjną lub generują nieakceptowalne opóźnienia w kontekście rozmowy głosowej \cite{Nguyen2025SBC}. Szczegółowe porównanie badanych modeli przedstawiono w sekcji~\ref{sec:modele}.

\subsection{Kwantyzacja dla inferencji na sprzęcie bez GPU}\label{sec:kwantyzacja}

Kwantyzacja jako technika redukcji precyzji numerycznej parametrów modelu znajduje zastosowanie w dwóch odrębnych kontekstach: optymalizacji procesu trenowania (omówionej w podsekcji~\ref{sec:finetuning}) oraz optymalizacji samej inferencji. Niniejsza podsekcja dotyczy wyłącznie drugiego zastosowania --- zmniejszenia precyzji wag gotowego modelu w celu przyspieszenia generacji tokenów i redukcji zużycia pamięci operacyjnej podczas działania.

Generatywna inferencja dużych modeli językowych jest procesem ograniczonym przepustowością pamięci (ang. \textit{memory-bandwidth-bound}), a nie wydajnością obliczeniową. Kim et al. \cite{Kim2023SqueezeLLM} wykazali, że generowanie każdego tokenu wymaga jednorazowego załadowania wszystkich wag z pamięci RAM do rejestrów procesora, wobec czego opóźnienie inferencji maleje w przybliżeniu liniowo wraz ze zmniejszaniem precyzji bitowej wag, niezależnie od narzutu związanego z dekwantyzacją. W środowisku CPU o architekturze ARM pozbawionym GPU, gdzie przepustowość pamięci wynosi ok. 50~GB/s (Raspberry Pi~5), model o 4~mld parametrów w precyzji FP16 zajmuje ok. 8~GB i musi być odczytany przy generowaniu każdego tokenu; po kwantyzacji 4-bitowej (Q4\_K\_M) zajmuje on ok. 2,5~GB, co przekłada się na ok. trzykrotnie wyższą przepustowość generacji.

Frantar et al. \cite{Frantar2023} zaproponowali metodę GPTQ (ang. \textit{Generative Pre-Trained Transformers Quantization}) --- jednorazową kwantyzację wag po zakończeniu trenowania (ang. \textit{post-training quantization}, PTQ) opartą na przybliżonych informacjach drugiego rzędu; pozwala ona skwantyzować model o 175~mld parametrów do 3--4 bitów w ok. cztery godziny pracy GPU, przy pomijalnym spadku jakości. GPTQ stanowi algorytmiczny fundament formatu GGUF stosowanego przez bibliotekę llama.cpp --- domyślny silnik inferencji w narzędziu Ollama, wykorzystanym w niniejszym projekcie.

Format GGUF (ang. \textit{GPT-Generated Unified Format}) przechowuje wagi w zgrupowanych schematach kwantyzacji (ang. \textit{groupwise quantization}), w których grupy 32 lub 256 wag współdzielą wspólny współczynnik skalowania. Poszczególne schematy różnią się precyzją: Q4\_K\_M (4 bity, wariant \textit{medium}) stanowi kompromis między rozmiarem a jakością; Q8\_0 (8 bitów) zapewnia wyższą precyzję kosztem większego zużycia pamięci; natomiast Q2\_K i Q3\_K agresywnie redukują rozmiar przy zauważalnym spadku jakości \cite{Kurt2026}. Schemat Q4\_K\_M stał się de facto standardem wdrożeń na architekturze ARM \cite{Renney2026}; w modelach typu mieszanki ekspertów (ang. \textit{Mixture-of-Experts}, MoE) agresywna kwantyzacja do Q2\_K może natomiast powodować kolaps jakości ze względu na wrażliwość mechanizmu routingu ekspertów na zakłócenia precyzji.

Implikacje opisanego mechanizmu dla wyników własnego benchmarku omówiono w sekcji~\ref{subsec:kwantyzacja}. Warto podkreślić, że schemat Q8\_0, wbrew intuicji, może okazać się wolniejszy od Q4\_K\_M na CPU pozbawionym akceleratora: wyższa precyzja oznacza większy transfer danych przy generowaniu każdego tokenu, co w środowisku ograniczonym przepustowością pamięci jest niekorzystne \cite{Tummalapalli2026}.

\subsection{Wielomodelowe systemy LLM — podział ról}

Pojęcie hierarchicznego podziału ról między modelami wywodzi się z destylacji wiedzy (ang. \emph{knowledge distillation}), wprowadzonej przez Hintona et al. \cite{Hinton2015}. Duży model nauczyciel (ang. \emph{teacher}) transferuje wiedzę do małego modelu ucznia (ang. \emph{student}) podczas trenowania, udostępniając miękkie etykiety (ang. \emph{soft labels}) zamiast twardych odpowiedzi. Podejście to zakłada wyraźną hierarchię zdolności: nauczyciel dysponuje bogatszą reprezentacją wiedzy, a uczeń odwzorowuje jego zachowanie przy istotnie mniejszych wymaganiach obliczeniowych. Autorzy wykazali, że tak przeniesiona wiedza pozwala kompaktowemu modelowi osiągnąć jakość zbliżoną do modelu źródłowego, co ustanowiło koncepcję współpracy modeli o różnej skali.

Koncepcja dynamicznego przydziału zadań rozwinęła się następnie w paradygmacie kaskadowego wnioskowania (ang. \emph{LLM cascading}). Chen et al. \cite{Chen2023FrugalGPT} wykazali, że sekwencyjne odpytywanie modeli w kolejności rosnącej zdolności — z zatrzymaniem po osiągnięciu progu jakości odpowiedzi — redukuje koszty inferencji nawet o 98\% przy jakości porównywalnej z korzystaniem wyłącznie z największego modelu. W tym ujęciu podział ról nie jest ustalony z góry, lecz wynika z trudności konkretnego zapytania: proste przypadki obsługuje tańszy model, a jedynie te wymagające trafiają do modelu o większych możliwościach.

Odmienny wzorzec współpracy stanowi dekodowanie spekulatywne (ang. \emph{speculative decoding}), zaproponowane przez Leviathana et al. \cite{Leviathan2023}. Mały i szybki model szkicowy (ang. \emph{draft model}) proponuje sekwencję kandydatów tokenów, które duży model docelowy weryfikuje w jednym przebiegu równoległym, akceptując lub odrzucając kolejne propozycje. Autorzy wykazali, że taka procedura przyspiesza generację bez zmiany rozkładu wyjściowego modelu docelowego. Podział ról wynika tu ze specjalizacji, a nie z trudności zadania: mały model odpowiada za tempo generacji, duży zaś gwarantuje ostateczną jakość.

Architektura Wilgi realizuje podział odmienny od wszystkich powyższych. Nauczyciel nie trenuje Ucznia w sposób permanentny, jak w destylacji wiedzy, nie weryfikuje jego wyjścia, jak w dekodowaniu spekulatywnym, ani nie przejmuje zapytań w sytuacji, gdy Uczeń nie radzi sobie z zadaniem, jak w kaskadzie modeli. Zamiast tego Nauczyciel wykonuje wyspecjalizowaną analizę danych sensorycznych IoT i produkuje ustrukturyzowany kontekst w postaci system promptu, natomiast Uczeń odpowiada wyłącznie za generację naturalnego języka polskiego na podstawie dostarczonego kontekstu. Podział kompetencji wynika zatem z charakteru zadania przypisanego każdemu z modeli, a nie z rozmiaru danych wejściowych ani ze złożoności pojedynczego zapytania.

\subsection{Inżynieria promptów i dostarczanie kontekstu}

Inżynieria promptów (ang. \emph{prompt engineering}) to zbiór technik kształtowania zapytań kierowanych do modelu językowego w celu poprawy jakości odpowiedzi oraz sterowania ich formą \cite{Schulhoff2024}. Schulhoff et al. dokonali systematycznej taksonomii ponad pięćdziesięciu technik promptowania, wyróżniając między innymi podejście bez przykładów (ang. \emph{zero-shot}), z kilkoma przykładami umieszczonymi w prompcie (ang. \emph{few-shot}) oraz z pośrednimi krokami rozumowania (ang. \emph{chain-of-thought}). Uporządkowanie tak licznych technik pokazuje, że sposób sformułowania zapytania stanowi odrębny wymiar sterowania modelem, niezależny od jego wag. Brown et al. \cite{Brown2020} wykazali, że duże modele adaptują zachowanie do zadania wyłącznie na podstawie przykładów zawartych w prompcie, bez modyfikacji wag — zjawisko to określa się jako uczenie w kontekście (ang. \emph{in-context learning}). Wei et al. \cite{Wei2022} pokazali z kolei, że wymuszenie pośrednich kroków rozumowania (ang. \emph{chain-of-thought}) znacząco poprawia jakość na zadaniach wieloetapowych, nawet bez dostrajania modelu. Obserwacje te dowodzą, że sama konstrukcja promptu może w istotnym stopniu zastąpić kosztowną modyfikację parametrów, co ma szczególne znaczenie w scenariuszach o ograniczonych zasobach.

W kontekście małych modeli językowych (ang. \emph{small language model}, SLM) działających na urządzeniach brzegowych szczególnego znaczenia nabiera ograniczanie zachowania modelu (ang. \emph{constraining}) za pomocą system promptu. Modele o rozmiarze rzędu 2--4 miliardów parametrów wykazują skłonność do generowania informacji nieobecnych w dostarczonym kontekście, czyli do halucynacji (ang. \emph{hallucination}), co przy braku dodatkowych zabezpieczeń podważa wiarygodność asystenta. Odpowiednio sformułowany system prompt może wymuszać określoną politykę zachowania — na przykład jawne komunikowanie braku danych, gdy wymagana informacja nie została przekazana — i tym samym ograniczać halucynacje bez modyfikacji wag modelu. Mechanizm ten jest tym cenniejszy, że nie wymaga ani dodatkowego treningu, ani zwiększenia rozmiaru modelu, a więc mieści się w ograniczeniach sprzętowych urządzenia brzegowego. Znaczenie tej techniki potwierdzają wyniki własnego benchmarku, omówione w sekcji~\ref{subsec:hallucination_dominance}.

Inspiracją dla sposobu dostarczania wiedzy zewnętrznej jest podejście Retrieval-Augmented Generation (RAG), zdefiniowane przez Lewisa et al. \cite{Lewis2020} jako połączenie parametrycznej pamięci modelu z nieparametryczną wiedzą zewnętrzną pozyskiwaną dynamicznie przez przeszukiwanie (ang. \emph{retrieval}). Klasyczny RAG pobiera dokumenty z bazy wektorowej na podstawie podobieństwa semantycznego i dołącza je do kontekstu zapytania. W Wildze dane zewnętrzne dostarczane są w odmienny sposób: Nauczyciel generuje deterministycznie ustrukturyzowany kontekst na podstawie aktualnych odczytów czujników IoT i przekazuje go Uczniowi jako system prompt. Nie jest przy tym wykonywana operacja wyszukiwania — kontekst jest każdorazowo konstruowany od podstaw z surowych danych. Takie podejście można określić mianem strukturyzowanego wstrzykiwania kontekstu (ang. \emph{structured context injection}), czerpiącego z filozofii RAG, czyli rozszerzania wiedzy modelu bez modyfikacji jego wag, lecz różniącego się mechanizmem dostarczania.

Opisane właściwości bezpośrednio kształtują projekt promptu Nauczyciela w Wildze. Prompt ten musi skłonić model do wytworzenia ustrukturyzowanego wyjścia, obejmującego pola analityczne takie jak stan energetyczny, komfort czy wzorce czasowe, zamiast swobodnej narracji, przy jednoczesnym zachowaniu rozumowania wykraczającego poza proste reguły deterministyczne. Jednocześnie prompt Ucznia musi wykorzystywać opisane techniki ograniczania zachowania, aby wygenerowana odpowiedź pozostała wierna dostarczonemu kontekstowi. Szczegółowy projekt tego promptu opisano w sekcji~\ref{subsec:przeplyw}.

\subsection{Optymalizacja hiperparametrów generacji}

Parametry generacji modelu językowego --- temperatura próbkowania (ang. \textit{sampling temperature}), rozmiar okna kontekstu (ang. \textit{context window}), maksymalna długość odpowiedzi, współczynnik top-p oraz kara za powtórzenia --- istotnie wpływają na jakość i szybkość generacji, a ich optymalne wartości zależą od modelu, zadania i sprzętu. Przeszukiwanie pełnej siatki kombinacji jest kosztowne obliczeniowo: przy pięciu parametrach i kilku wartościach każdego liczba kombinacji szybko przekracza możliwości wyczerpującej eksploracji.

Akiba et al. \cite{Akiba2019} zaproponowali bibliotekę Optuna --- framework optymalizacji hiperparametrów oparty na bayesowskim próbkowaniu z algorytmem TPE (ang. \textit{Tree-structured Parzen Estimator}). Zamiast testować kombinacje losowo lub w siatce, TPE buduje probabilistyczny model funkcji celu na podstawie dotychczasowych wyników i kieruje kolejne próby ku obiecującym obszarom przestrzeni parametrów, osiągając porównywalną eksplorację przy znacznie mniejszej liczbie ewaluacji. W benchmarku zastosowano 30 triali Optuny na model (co przy pełnej siatce odpowiadałoby 360 kombinacjom), uzyskując dodatkowo histogram ważności cech (ang. \textit{feature importance}), który sam stanowi wynik analityczny: pokazuje, które parametry generacji mają istotny wpływ na composite score.

\subsection{Dostrajanie modeli językowych}\label{sec:finetuning}

O ile kwantyzacja omówiona w podsekcji~\ref{sec:kwantyzacja} dotyczy redukcji precyzji wag gotowego modelu wyłącznie w celu optymalizacji inferencji, dostrajanie (ang. \textit{fine-tuning}) oznacza modyfikację tych wag --- adaptację modelu do nowych zadań lub domeny poprzez aktualizację parametrów na zbiorze treningowym.

Pełne dostrajanie (ang. \textit{full fine-tuning}) modeli o miliardach parametrów jest niepraktyczne na sprzęcie konsumenckim. Hu et al. \cite{Hu2022} zaproponowali metodę LoRA (ang. \textit{Low-Rank Adaptation}): zamiast modyfikować wszystkie wagi, dodaje się pary małych macierzy o niskim rzędzie (ang. \textit{low-rank matrices}), których iloczyn aproksymuje zmianę wag, podczas gdy wagi bazowe pozostają zamrożone. LoRA redukuje liczbę trenowalnych parametrów nawet 10~000-krotnie (dla modelu GPT-3 o 175~mld parametrów), przy jakości porównywalnej lub wyższej od pełnego dostrajania na standardowych zadaniach przetwarzania języka naturalnego (ang. \textit{NLP}); dodatkowe macierze można scalić z wagami bazowymi po zakończeniu trenowania, eliminując narzut podczas inferencji.

Dettmers et al. \cite{Dettmers2023} rozszerzyli metodę LoRA o kwantyzację wag modelu bazowego podczas trenowania. W metodzie QLoRA wagi zamrożonego modelu bazowego przechowywane są w precyzji 4-bitowej (NF4), natomiast gradienty obliczane są w wyższej precyzji dla zachowania stabilności numerycznej; umożliwia to dostrajanie modeli o 7~mld parametrów na GPU z 24~GB pamięci VRAM, a nawet modeli o 65~mld parametrów na 48~GB VRAM. W kontekście niniejszej pracy QLoRA stanowi metodologiczną podstawę planowanego eksperymentu dostrajania Ucznia na polskojęzycznych danych z domeny inteligentnego domu (ang. \textit{smart-home}), z wykorzystaniem biblioteki mlx-lm na platformie Apple M4.

Planowany eksperyment ma odpowiedzieć na pytanie, czy model dostrojony na 400 przykładach --- bez jawnych reguł zawartych w promptcie systemowym --- osiąga wyższą jakość odpowiedzi niż model bazowy z precyzyjnie zdefiniowanymi regułami w promptcie systemowym. Wynik dostarczyłby empirycznej podstawy do oceny, w jakim stopniu dostrajanie może zastąpić inżynierię promptów jako mechanizm kontroli zachowania małego modelu językowego (ang. \textit{SLM}).

\subsection{Interfejsy głosowe: rozpoznawanie mowy, detekcja aktywności głosowej i synteza mowy}

Kompletny pipeline głosowy wymaga mechanizmu detekcji aktywności głosowej (ang. \textit{Voice Activity Detection}, VAD), odróżniającego segmenty mowy od ciszy i szumów tła. Bez VAD model rozpoznawania mowy przetwarza ciągły strumień audio, generując fałszywe transkrypcje z segmentów ciszy i zwiększając latency. Silero Team \cite{SileroVAD} opracowali lekki model neuronowy VAD do wdrożeń brzegowych: przetwarza fragmenty audio o długości około 30 ms w czasie poniżej 1 ms na CPU, przy rozmiarze modelu poniżej 1 MB; wspiera próbkowanie 8 kHz i 16 kHz, dostępny jest na licencji MIT oraz zoptymalizowany dla PyTorch i ONNX Runtime.

W rozpoznawaniu mowy (ang. \textit{Speech-to-Text}, STT) modelem referencyjnym w open-source jest Whisper, opracowany przez Radford et al. \cite{Radford2023}. Trenowano go metodą słabego nadzoru (ang. \textit{weak supervision}) na 680\,000 godzin wielojęzycznego audio, co daje wysoką odporność na szumy tła --- kluczową w warunkach domowych o nieregularnym otoczeniu akustycznym. Whisper wspiera transkrypcję polskiego i jest dostępny w kilku wariantach rozmiaru; wariant \textit{faster-whisper} (implementacja oparta na CTranslate2) stosowany jest w projekcie ze względu na niższe wymagania pamięciowe i wyższą przepustowość na CPU.

Dla syntezy mowy (ang. \textit{Text-to-Speech}, TTS) modelem referencyjnym jest VITS \cite{Kim2021}: warunkowy autoenkoder wariacyjny połączony z trenowaniem adversarialnym, umożliwiający generację naturalnie brzmiącego głosu metodą end-to-end przy niskim latency. Komponent TTS zaplanowano jako element docelowej architektury Wilgi; w obecnym prototypie nie został jeszcze zaimplementowany.
